% !TEX program = pdflatex
\documentclass[a4paper,12pt]{extreport}
\usepackage[T5]{fontenc}
\usepackage[utf8]{inputenc}
\usepackage{mathptmx}
\usepackage[vietnamese]{babel}
\usepackage[left=3.5cm,right=2cm,top=2cm,bottom=2cm]{geometry}
\usepackage{setspace}\setstretch{1.2}
\usepackage{amsmath,amssymb}
\usepackage{graphicx,float}
\usepackage{listings}
\usepackage[hidelinks]{hyperref}

\begin{document}

\chapter*{BÀI MÔ TẢ HỆ THỐNG CHATBOT AI TÍCH HỢP TRÍCH XUẤT TÀI LIỆU}

Hệ thống là giải pháp kết hợp giữa trích xuất tài liệu tự động và chatbot
trí tuệ nhân tạo, cho phép người dùng đặt câu hỏi về nội dung tài liệu và
nhận câu trả lời chính xác kèm trích dẫn nguồn. Hệ thống gồm ba tầng:
tầng trích xuất xử lý tài liệu đầu vào thành dữ liệu có cấu trúc, tầng
lưu trữ và truy xuất quản lý dữ liệu đã trích xuất, và tầng chatbot trả
lời câu hỏi bằng mô hình ngôn ngữ lớn. Kiến trúc này cho phép hệ thống
áp dụng cho nhiều lĩnh vực: từ tra cứu pháp luật giao thông, quản lý nhà
trọ, đến xử lý chứng từ tài chính và hồ sơ nghiệp vụ.


\section*{Tầng trích xuất tài liệu}

Tầng trích xuất tiếp nhận tài liệu đầu vào ở nhiều định dạng: PDF, ảnh
scan, bảng tính. Hệ thống sử dụng Docling để phân tích cấu trúc tài liệu,
nhận diện bảng biểu, hình ảnh và luồng văn bản. Kết quả phân tích được
ánh xạ vào mô hình tài liệu chuẩn (Canonical Document Model, CDM) với
cấu trúc phân cấp: tài liệu chứa các phần, mỗi phần chứa các đoạn văn
bản, bảng biểu và hình ảnh kèm thông tin vị trí trang và dòng.

Sau khi có CDM, tầng xử lý quy tắc (Rule Engine) thực hiện trích xuất
các trường thông tin quan trọng dựa trên tệp cấu hình theo lĩnh vực. Quy
trình gồm bốn bước: xác định vị trí thông tin trong tài liệu (locate),
trích xuất giá trị (extract), chuẩn hóa dữ liệu (normalize) và kiểm tra
tính hợp lệ (validate). Mỗi giá trị trích xuất đều kèm dấu vết nguồn
(provenance) ghi lại trang và dòng trong tài liệu gốc, phục vụ kiểm toán
và đối chiếu sau này.

Hệ thống hỗ trợ kiến trúc plugin cho phép mở rộng sang lĩnh vực mới mà
không cần sửa đổi phần lõi. Nhà phát triển chỉ cần viết plugin và tệp
cấu hình quy tắc cho lĩnh vực đó; phần xử lý tài liệu, nhận diện chữ
viết, phân tích bảng biểu và kiểm tra dữ liệu do hệ thống đảm nhiệm.


\section*{Tầng lưu trữ và truy xuất}

Dữ liệu sau trích xuất được tổ chức thành các đoạn văn bản (chunk) theo
cấu trúc logic của tài liệu gốc. Mỗi chunk mang metadata giàu: định danh
tài liệu, vị trí trong cấu trúc phân cấp, và dấu vết nguồn. Các chunk
được nhúng thành vector 384 chiều bằng mô hình
intfloat/multilingual-e5-small và lưu vào cơ sở dữ liệu vector ChromaDB.

Hệ thống sử dụng truy xuất lai (hybrid retrieval) kết hợp hai phương
pháp: tìm kiếm ngữ nghĩa bằng vector embedding (dense retrieval) và tìm
kiếm từ khóa BM25 Okapi (sparse retrieval). Kết quả từ hai phương pháp
được hợp nhất bằng thuật toán Reciprocal Rank Fusion (RRF). Metadata từ
tầng trích xuất cho phép lọc chính xác theo lĩnh vực, loại tài liệu,
hoặc khoảng thời gian.


\section*{Tầng chatbot trả lời câu hỏi}

Khi người dùng gửi câu hỏi, tầng chatbot kích hoạt chuỗi xử lý trạng
thái sử dụng LangGraph. Quy trình gồm các bước:

\begin{itemize}
    \item Viết lại câu hỏi để tối ưu truy xuất.
    \item Truy xuất lai trên ChromaDB và BM25 để lấy các chunk liên quan.
    \item Chấm điểm chất lượng kết quả truy xuất bằng mô hình ngôn ngữ.
    \item Nếu điểm thấp, viết lại câu hỏi và truy xuất lại hoặc chuyển
          sang tìm kiếm web dự phòng.
    \item Sinh câu trả lời dựa trên ngữ cảnh các chunk đã truy xuất.
    \item Kiểm tra từng trích dẫn trong câu trả lời có khớp với nguồn
          thực tế hay không.
    \item Stream câu trả lời đã kiểm duyệt về giao diện qua SSE.
\end{itemize}

Mô hình ngôn ngữ chính là Google Gemini 2.5 Flash với free tier 250
request mỗi ngày. Khi vượt định mức hoặc gặp lỗi, hệ thống tự động
chuyển sang OpenAI GPT-4o-mini. Mô hình phụ này cũng đóng vai trò giám
khảo độc lập trong quy trình đánh giá chất lượng, tránh thiên vị do cùng
một mô hình vừa sinh câu trả lời vừa chấm điểm.

Câu trả lời trả về dưới dạng Markdown kèm danh sách trích dẫn có cấu
trúc: tên văn bản, số Điều, Khoản, Điểm, trạng thái hiệu lực và dấu vết
nguồn (trang, dòng). Mọi câu trả lời đều kèm điều khoản miễn trừ trách
nhiệm pháp lý.


\section*{Vai trò người dùng}

\subsection*{Người dùng cuối}

Người dùng cuối tương tác với hệ thống qua giao diện trò chuyện. Họ nhập
câu hỏi bằng ngôn ngữ tự nhiên và nhận câu trả lời ngắn gọn kèm trích
dẫn. Hệ thống phục vụ nhiều đối tượng: người tham gia giao thông tra cứu
mức phạt, sinh viên luật tìm điều khoản, chủ nhà trọ tra cứu hợp đồng
thuê, nhân viên tín dụng kiểm tra lịch trả nợ.

\subsection*{Nhà quản trị và phát triển}

Nhà quản trị chịu trách nhiệm cập nhật cơ sở dữ liệu tài liệu và theo
dõi hiệu năng hệ thống. Họ nạp tài liệu mới qua API được xác thực bằng
Bearer Token; quy trình tự động thực hiện trích xuất, nhúng vector và tái
xây dựng chỉ mục BM25 mà không làm gián đoạn luồng tra cứu. Nhà phát
triển mở rộng hệ thống sang lĩnh vực mới bằng cách viết plugin và tệp
cấu hình quy tắc, không cần can thiệp vào phần lõi xử lý.


\section*{Kiến trúc hệ thống}

Hệ thống vận hành trên kiến trúc đa tầng:

\begin{itemize}
    \item FastAPI làm khung API cho cả tầng trích xuất và chatbot.
    \item LangGraph làm bộ điều phối trạng thái cho luồng xử lý chatbot.
    \item ChromaDB làm cơ sở dữ liệu vector cho truy xuất ngữ nghĩa.
    \item SQLite cho lưu trữ trạng thái hội thoại và lịch sử.
    \item Hệ thống tệp cho tài liệu PDF gốc và chỉ mục BM25.
    \item Công cụ dòng lệnh cho quy trình nạp dữ liệu hàng loạt.
\end{itemize}

Hệ thống chạy trên một worker duy nhất do ChromaDB dùng SQLite nội bộ,
không hỗ trợ nhiều tiến trình ghi đồng thời. Toàn bộ mô hình ngôn ngữ và
nhúng vector chạy qua API cloud; card đồ họa MX330 2GB VRAM trên máy
phát triển không đủ để chạy inference cục bộ.


\section*{Bảo mật}

Hệ thống tuân thủ Nghị định 13/2023/NĐ-CP về bảo vệ dữ liệu cá nhân:
không lưu trữ thông tin nhận dạng ngoài mã định danh phiên, chính sách
lưu trữ tối đa 30 ngày, xác thực API bằng Bearer Token, giới hạn tốc độ
6 request mỗi phút mỗi địa chỉ IP.

\end{document}
